\def\topic{Capture-avoiding de Bruijn substitution}
\ifdefined\MyHeaderIncluded
\else
\def\qwertyflag{true}

\ifdefined\MyHeaderIncluded

\else
  \newcommand{\foobar}{FooBar}
  \def\MyHeaderIncluded{true}
% if is till the end of the file

\documentclass[10pt,a4paper]{article}
\usepackage[utf8]{inputenc}
%\usepackage[german]{babel}
%\usepackage[cyrillic]{babel}
\usepackage[T1]{fontenc}
%\usepackage[T2A]{fontenc}
%\usepackage[english, russian]{babel}
\usepackage{amsmath}
\usepackage{amsfonts}
\usepackage{amssymb}
\usepackage{amsthm}
\usepackage{mathtools} %coloneqq

\usepackage{bussproofs} %

%\renewcommand{\thesubsection}{\arabic{subsection}}
%\newtheorem{thm}{Theorem}[section]
%\renewcommand{\thethm}{\arabic{subsection}.\arabic{thm}}

\usepackage{thmtools}
\usepackage{hyperref}
\hypersetup{
    colorlinks=true,
    linkcolor=blue,
    filecolor=magenta,      
    urlcolor=cyan,
    pdftitle={Overleaf Example},
    %pdfpagemode=FullScreen,
    }
%\usepackage{stix-type1}%\turnediota
\theoremstyle{definition}

\newtheorem{constructor}{Constructor}[section]
\newtheorem{iaxiom}{Axiom}
\newtheorem{axiom}{Axiom}
\newtheorem*{axiom*}{Axiom}
\newtheorem{definition}{Definition}[section]
\newtheorem*{definition*}{Definition}
\newtheorem{remark}{Remark}[section]
\newtheorem*{remark*}{Remark}
\newtheorem{recall}{Recall}[section]
\newtheorem*{recall*}{Recall}

\newtheorem{lemma}{Lemma}[section]
\newtheorem{theorem}{Theorem}[section]
\newtheorem*{theorem*}{Theorem}
\newtheorem{note}{Note}[section]
\usepackage{etoolbox}

%\newcommand*{\@Dashv}[2]{%
%  \reflectbox{$\m@th#1#2$}%
%}
\newcommand*{\separation}[3]{% 1)binder, 2)set, 3)condition
  \{#1\in #2: #3\}
}

% https://www.freecodecamp.org/news/simple-regex-tricks-for-beginners-3acb3fa257cb/
%  \\mathcal\{T\} |---> \topology
%  \\mathcal\{([^\}]*)\}   |--->  $1
\newcommand{\AND}{\ \ \&\ \ }
\newcommand{\any}{\mathrm{\ any}}
\newcommand{\borel}{\mathcal{B}}
\newcommand{\can}{\mathrm{can}}
\newcommand{\canonize}{\mathrm{canonize}}
\newcommand{\closure}{\mbox{cl}}
\newcommand{\complex}{\mathbb{C}}
\newcommand{\converges}{\longrightarrow}
\newcommand{\countable}{\mbox{\ countable}}
\newcommand{\ldefined}{\mbox{\ defined}}
\newcommand{\dom}{\mbox{dom}}
\newcommand{\Dom}{\mbox{Dom}}
\newcommand{\Done}{\noindent \textbf{Done.}\\}
\newcommand{\lexists}{\mbox{\ exists}}
\newcommand{\lforall}{\mbox{\ forall}}
\newcommand{\lunique}{\mbox{\ unique}}
\newcommand{\family}{\mathrm{\ family}}
\newcommand{\Fm}{\mathrm{Fm}}
\newcommand{\FuncSymb}{\mathrm{FuncSymb}}
\newcommand{\function}{\mathrm{\ function}}
\newcommand{\functional}{\mbox{\ functional}}
\newcommand{\IFF}{\ \Longleftrightarrow\ }
\renewcommand{\iff}{\ \longleftrightarrow\ }
\newcommand{\IMPLIES}{\ \Longrightarrow\ }
\renewcommand{\implies}{\ \longrightarrow\ }
\newcommand{\Inf}{\mathrm{Inf}}
\newcommand{\injective}{\mbox{\ injective}}
\newcommand{\injection}{\mbox{\ injection}}
\newcommand{\interior}{\mbox{int}}
\newcommand{\intersection}{\bigcap}
\newcommand{\isempty}{\mathrm{\ empty}}
\newcommand{\LB}{\mathrm{LB}}
\newcommand{\nghbset}{\mathcal{N}}
\newcommand{\nonempty}{\mathrm{\ nonempty}}
\newcommand{\Ord}{\mathrm{Ord}}
\newcommand{\power}{\mathcal{P}}
\newcommand{\pr}{\mathrm{pr}}
\newcommand{\fst}{\mathrm{fst}}
\newcommand{\snd}{\mathrm{snd}}
\newcommand{\PredSymb}{\mathrm{PredSymb}}
\newcommand{\Proof}{\noindent \textbf{Proof.}\\}
\newcommand{\ran}{\mathrm{ran}}
\newcommand{\Ran}{\mathrm{Ran}}
\newcommand{\real}{\mathbb{R}}
\newcommand{\relation}{\mathrm{\ relation}}
\newcommand{\set}{\mathrm{\ set}}
\newcommand{\singleton}{\mathrm{\ singleton}}
\newcommand{\semantics}{\mathrm{semantics}}
\newcommand{\semialgebra}{\mathcal{S}}
\newcommand{\Sup}{\mathrm{Sup}}
\newcommand{\surjective}{\mbox{\ surjective\ on\ }}
\newcommand{\Term}{\mathrm{Term}}
\newcommand{\Type}{\mathrm{Type}}
\newcommand{\topology}{\mathcal{T}}
\newcommand{\total}{\mbox{\ total\ on\ }}
\newcommand{\transitive}{\mathrm{\ transitive}}
\newcommand{\union}{\bigcup}
\newcommand{\lTrue}{\mathrm{True}}
\newcommand{\lFalse}{\mathrm{False}}
\newcommand{\lUncertain}{\mathrm{Uncertain}}
\newcommand{\lif}{\mathrm{if}}
\newcommand{\lthen}{\mathrm{then}}
\newcommand{\lelse}{\mathrm{else}}
\newcommand{\lexception}{\mathrm{exception}}
\newcommand{\lundef}{\mathrm{undef}}
\newcommand{\univocal}{\mathrm{\ univocal}}
\newcommand{\univocalat}{\mathrm{\ univocal\ at\ }}
\newcommand{\UB}{\mathrm{UB}}
\newcommand{\Un}{\mathcal{U}}
\newcommand{\Var}{\mathrm{Var}}

\newcommand{\notvarphi}{\phi}
\renewcommand{\phi}{\varphi}

\newcommand{\comment}[1]{\quad\mbox{(#1)}}
% local commands
%\preto\section{\undef\blipo} %etoolbox
\preto\subsection{\undef\blipo} %etoolbox
% https://tex.stackexchange.com/questions/105684/limit-macro-newcommand-scope-to-section
\usepackage{graphicx}
\usepackage{authblk}

\setlength\parindent{0pt} % whole doc of \noindent
\usepackage{soul}
%\usepackage{mathgl}

\fi % guard
\ifdefined\topic
\else
\def\topic{No Titel}
\fi
\title{\topic}
\author[1]{Georgy Dunaev}
\begin{document}
\maketitle
\fi
\section{A complete de Bruijn substitution example}
\label{sec:debruijn_substitution}
The preceding listing is preserved as a historical prototype. The
following construction completes its substitution and stopping cases
in a separate, executable example. Terms here are untyped finite trees
\[
t::=\operatorname{Var}(k)\mid\operatorname{Lam}(t)
       \mid\operatorname{App}(t,t),\qquad k\in\mathbb N.
\]
Index zero refers to the nearest binder or the first slot in the
surrounding context. Open terms are allowed. This implementation is
not an implementation of the typed calculus or checker developed in
the next section, and no normalization claim for arbitrary untyped
terms is made.

\subsection{Scope and insertion of context slots}
Write $\operatorname{sc}_n(t)$ when $t$ is well scoped in a context
of $n$ slots, numbered $0,\ldots,n-1$ from nearest to furthest:
\[
\begin{aligned}
\operatorname{sc}_n(\operatorname{Var}(k))&\iff 0\leq k<n,\\
\operatorname{sc}_n(\operatorname{Lam}(b))&\iff\operatorname{sc}_{n+1}(b),\\
\operatorname{sc}_n(\operatorname{App}(f,a))
 &\iff\operatorname{sc}_n(f)\land\operatorname{sc}_n(a).
\end{aligned}
\]
In particular $\operatorname{sc}_0$ describes closed terms.
For nonnegative integers $a,c$, define a shift $\uparrow^a_c$ by
\[
\begin{aligned}
\uparrow^a_c\operatorname{Var}(k)
 &=\begin{cases}\operatorname{Var}(k),&k<c,\\
                  \operatorname{Var}(k+a),&k\geq c,\end{cases}\\
\uparrow^a_c\operatorname{Lam}(b)
 &=\operatorname{Lam}(\uparrow^a_{c+1}b),\\
\uparrow^a_c\operatorname{App}(f,u)
 &=\operatorname{App}(\uparrow^a_c f,\uparrow^a_c u).
\end{aligned}
\]
The cutoff protects the nearest $c$ slots; it increases under a binder
so that variables bound inside the term are protected as well.

\begin{lemma}[Insertion of slots preserves scope and variable reference]
\label{lem:debruijn_shift_scope}
If $\operatorname{sc}_{n+c}(t)$, then
$\operatorname{sc}_{n+c+a}(\uparrow^a_c t)$.
More precisely, let $\Delta$ be a list of $c$ distinct context names
and $\Gamma$ a list of $n$ further names. Inserting $a$ fresh names
between $\Delta$ and $\Gamma$ while shifting by $\uparrow^a_c$
preserves the named term represented by $t$, up to renaming its bound
variables.
\end{lemma}
\Proof
Induct on $t$. A variable $k<c$ remains in the same slot of $\Delta$.
For $k\geq c$, its new index $k+a$ refers to the same slot of
$\Gamma$ and is less than $n+c+a$. An abstraction puts one new name
at the front of $\Delta$; apply the induction hypothesis with $c+1$
to its body and retain the abstraction node. For an application apply
the hypothesis to both children. Named binders may be chosen fresh
from an infinite set of names; all terms and context lists involved
are finite. Thus no inserted name is captured, proving both assertions.
\Done

\subsection{Removing one binder by substitution}
Let $R_d(s,t)$ remove slot $d$ from $t$ and replace its occurrences
by the outer-context term $s$. The nearest $d$ slots are local binders
crossed on the way to the occurrence. Define
\[
\begin{aligned}
R_d(s,\operatorname{Var}(k))
 &=\begin{cases}
   \operatorname{Var}(k),&k<d,\\
   \uparrow^d_0s,&k=d,\\
   \operatorname{Var}(k-1),&k>d,
   \end{cases}\\
R_d(s,\operatorname{Lam}(b))
 &=\operatorname{Lam}(R_{d+1}(s,b)),\\
R_d(s,\operatorname{App}(f,u))
 &=\operatorname{App}(R_d(s,f),R_d(s,u)).
\end{aligned}
\]
The argument $s$ remains in its original outer context during the
recursion. It is shifted only when inserted. The final variable case
decrements indices because a context slot has been removed, and
cannot produce a negative index since $k>d\geq0$.

\begin{theorem}[Scope and capture avoidance for binder removal]
\label{thm:debruijn_replacement}
If $\operatorname{sc}_n(s)$ and $\operatorname{sc}_{n+d+1}(t)$, then
$\operatorname{sc}_{n+d}(R_d(s,t))$.
Read $t$ in the named context $\Delta,x,\Gamma$, where
$|\Delta|=d$, $|\Gamma|=n$, and all names are distinct, and read $s$
in $\Gamma$. Then reading $R_d(s,t)$ in $\Delta,\Gamma$ gives the
capture-avoiding substitution of the reading of $s$ for $x$ in the
reading of $t$, up to bound-variable renaming.
\end{theorem}
\Proof
Induct on $t$, allowing every $d$ in the induction hypothesis.
If $t=\operatorname{Var}(k)$ and $k<d$, its unchanged index is less
than $n+d$ and names the same local variable. If $k=d$, the shift
lemma with cutoff zero gives
$\operatorname{sc}_{n+d}(\uparrow^d_0s)$; its named reading is precisely
that of $s$, with the $d$ local names inserted before its context.
Those names are distinct from $\Gamma$, so this insertion captures no
free variable of $s$. If $k>d$, the bound $k<n+d+1$ gives
$k-1<n+d$, and decrementing the index refers to the same name in
$\Gamma$ after $x$ has been removed.

For an abstraction, the body has scope $n+d+2$. Apply the induction
hypothesis with depth $d+1$, obtaining a body of scope $n+d+1$,
and restore the abstraction. For its named interpretation choose a
binder fresh for all context names and the free names of $s$.
The induction hypothesis therefore computes ordinary capture-avoiding
substitution under that binder. For an application, both scope and
the substitution identity follow from the two induction hypotheses.
These cases prove the result.
\Done

In particular, if $\operatorname{sc}_{n+1}(b)$ and
$\operatorname{sc}_n(s)$, beta contraction is
\[
\operatorname{App}(\operatorname{Lam}(b),s)
\longrightarrow R_0(s,b).
\]
The theorem proves both its scope preservation and agreement with the
named beta rule. This is binder-eliminating substitution, not an
operation which replaces a slot while keeping the old context length.
For example, in a one-slot outer context,
\[
\operatorname{App}(\operatorname{Lam}(\operatorname{Lam}(\operatorname{Var}(1))),
                   \operatorname{Var}(0))
\longrightarrow\operatorname{Lam}(\operatorname{Var}(1)).
\]
The outer variable remains free; returning
$\operatorname{Lam}(\operatorname{Var}(0))$ would capture it.

\subsection{One step, normal forms, and a finite computation budget}
The function \texttt{step} below contracts an outer application redex
first. Otherwise it searches an abstraction's body, or an application's
left child before its right child. Thus it performs one leftmost-outermost
compatible beta step, including steps under abstractions.

\begin{proposition}[The step and fuel contracts]
\label{prop:debruijn_step_fuel}
On valid finite terms, \texttt{step} terminates. It returns
\texttt{NONE} exactly for terms containing no beta redex; a result
\texttt{SOME u} is one compatible beta step to $u$. Every step
preserves $\operatorname{sc}_n$ for each applicable $n$.
For a nonnegative budget $N$, \texttt{run N t} terminates. A result
\texttt{Normal u} gives a beta-normal form reached in at most $N$
steps. A result \texttt{OutOfFuel u} means that exactly $N$ steps
were performed and $u$ still has a redex.
\end{proposition}
\Proof
Every recursive search call in \texttt{step} uses a proper subtree;
replacement and shifting are also structural recursions. Its variable
case has no redex, its abstraction case delegates exactly to the body,
and its application cases either contract the root redex or search
both children in the stated order. Induction on the tree therefore
proves termination, the normal-form test, and the single-step claim.
At a contracted root the replacement theorem preserves scope. Under
an abstraction use that theorem or the induction hypothesis in scope
$n+1$; under an application use scope $n$ for each child.

In \texttt{run}, a successful step either decreases the budget by one
or, if the budget is zero, returns the current term as
\texttt{OutOfFuel}. The loop therefore terminates. Induction on the
budget proves the two return-value assertions. At budget zero it
still tests for a redex, so an already normal term is reported as
normal, including one reached by the last allowed step. This test
computes a possible next term but does not commit an additional step.
\Done

Fuel exhaustion alone does not distinguish a divergent term from one
which needs more steps. There is, however, a direct divergence proof
for the historical example. Put
\[
\delta=\operatorname{Lam}(\operatorname{App}(\operatorname{Var}(0),
                                            \operatorname{Var}(0))),
\qquad \Omega=\operatorname{App}(\delta,\delta).
\]
Its only redex is the root, and contracting that redex returns
$\Omega$. Thus every finite reduction sequence stays at this nonnormal
term, so $\Omega$ has no beta-normal form, and \texttt{run} exhausts
every finite budget on it. On the other hand,
$\operatorname{App}(\operatorname{Lam}(I),\Omega)$, where
$I=\operatorname{Lam}(\operatorname{Var}(0))$, reaches $I$ in one
step: the argument need not be normalized before substitution.
No general theorem about the termination or completeness of unbounded
normal-order evaluation is needed for these assertions.

\subsection{Executable Standard ML}
The following file is self-contained. It uses \texttt{IntInf.int}
for indices and fuel. Its computational operations reject negative
indices, depths, cutoffs, shifts, or fuel with \texttt{Domain};
\texttt{scoped} instead returns false for a negative context size or
any negative index. The structural scope predicate remains an explicit test: a nonnegative
index alone need not fit a particular surrounding context. The public
signature hides the unchecked recursive helpers.
\begin{nolinenumbers}
\lstinputlisting{examples/debruijn.ML}
\end{nolinenumbers}
For example, after loading the file in Poly/ML, evaluate
\begin{lstlisting}
open DeBruijn;
val identity = Lam (Var 0);
val duplicate = Lam (App (Var 0, Var 0));
val omega = App (duplicate, duplicate);
val terminating = run 1 (App (Lam identity, omega));
val exhausted = run 20 omega;
\end{lstlisting}
The first result is \texttt{Normal identity}; the second is
\texttt{OutOfFuel omega}. Repository tests also compare bounded
families of open terms with an independent named-substitution oracle,
check scope preservation, and exercise the neutral-application stopping
case missing from the historical sketch.
\ifdefined\qwertyflag
\undef\qwertyflag
\end{document}
\fi
